% 严格保留原图左侧单元自循环与右侧按时间展开的四个单元。
\resizebox{\linewidth}{!}{%
\begin{tikzpicture}[
    font=\large,
    unit/.style={draw,line width=.9pt,fill=green!15,minimum width=1.15cm,minimum height=.75cm},
    inp/.style={circle,draw,line width=.8pt,fill=blue!25,minimum size=.66cm,inner sep=1pt},
    hidden/.style={inp,fill=violet!30},
    flow/.style={-{Stealth[length=2.5mm]},line width=1pt}
]
\node[unit] (a) at (0,0) {$A$};
\node[inp] (x) at (0,-1.45) {$x_t$};
\node[hidden] (h) at (0,1.45) {$h_t$};
\draw[flow] (x) -- (a);
\draw[flow,rounded corners=3pt] (a.east) -- (1.05,0) -- (1.05,.65) -- (-1.05,.65) -- (-1.05,0) -- (a.west);
\draw[flow,preaction={draw=white,line width=3pt}] (a) -- (h);
\node at (2.05,0) {$=$};
\foreach \id/\xx/\tt in {a0/3.7/0,a1/5.8/1,a2/7.9/2,at/11.1/t}{
  \node[unit] (\id) at (\xx,0) {$A$};
  \node[inp] (x\id) at (\xx,-1.45) {$x_{\tt}$};
  \node[hidden] (h\id) at (\xx,1.45) {$h_{\tt}$};
  \draw[flow] (x\id) -- (\id);
  \draw[flow] (\id) -- (h\id);
}
\draw[flow] (a0) -- (a1);
\draw[flow] (a1) -- (a2);
\draw[flow] (a2) -- (at);
\node at (9.45,-1.45) {$\cdots$};
\end{tikzpicture}%
}
